\section{Introduction}
A widespread concern with large language models is hallucination: the
production of fluent text that is not factually grounded. Recent work
has approached the problem from the activation side
\cite{azaria2023internal,burns2023discovering}; in this paper we ask
whether a purely structural signal---the topology of the per-layer
attention graph---suffices.

\paragraph{Question.} \emph{Do truthful and hallucinated continuations
yield statistically distinguishable persistent-homology summaries of
the attention graph?}

\paragraph{Setup.} For each factual prompt we generate two
continuations: one factually correct (e.g.\ ``Paris'') and one
plausible but wrong (``Berlin''). The surface form is matched. We
forward both through GPT-2, extract every layer's attention matrix,
convert to a Vietoris--Rips filtration over the metric
$d_{ij}=1-\max(A_{ij},A_{ji})$, and compute persistence in $H_0$ and
$H_1$.

\paragraph{Pilot.} The local configuration uses $\gptTwoNumPairs$ curated
pairs; this is intentionally underpowered. The point of the pilot is to
(i) validate the pipeline and (ii) measure the effect size so that the
GPU configuration can be sized correctly. The local pilot finds no
layer with $p<0.05$ on $H_0$ (minimum $p=\gptTwoMinPval$ at layer
$\gptTwoMinPvalLayer$); under the observed effect-size estimate, a
sample of $N\ge 200$ pairs would be required to reach $\alpha=0.05$ at
$\beta=0.2$, which is exactly what the
\texttt{gpu\_full.yaml} configuration supplies via TruthfulQA.

\paragraph{Contributions.} (i) An open, reproducible pipeline for
attention-graph persistent homology; (ii) a curated factual pair
dataset (10 pairs in this release); (iii) effect-size estimation for
$H_0$ and $H_1$ persistence to plan future runs.
